%% OPERA: Robust Representation Learning for Irregular Time Series under Observation-Process Shift
%% Draft 2026-09-24. All reported numbers are single-seed (seed 0) and marked with \prelim; the protocol
%% requires 5 seeds before any claim is final. Do not remove a \prelim or a \todo without replacing the
%% underlying measurement.
\documentclass[sigconf,nonacm]{acmart}
\usepackage{booktabs}
\usepackage{amsmath}
\usepackage{amssymb}
\usepackage{algorithm}
\usepackage{algpseudocode}
\usepackage{xcolor}
\newcommand{\todo}[1]{\textcolor{red}{[#1]}}
\newcommand{\prelim}{\textsuperscript{\dag}}
\newcommand{\Sobs}{\mathcal{S}}
\newcommand{\Oobs}{\mathcal{O}}

\AtBeginDocument{\providecommand\BibTeX{{Bib\TeX}}}
\settopmatter{printacmref=false}
\renewcommand\footnotetextcopyrightpermission[1]{}
\pagestyle{plain}

\begin{document}

\title{OPERA: Robust Representation Learning for Irregular Time Series\\ under Observation-Process Shift}

\author{Anonymous Authors}
\affiliation{\institution{Anonymous Institution}\country{}}

\begin{abstract}
An irregularly sampled time series is not a signal but the record of a measurement procedure: an underlying
temporal process observed through a process that decides when, and which variables, get recorded. Methods for
such data are evaluated almost exclusively under the observation process of the benchmark they are trained on.
We ask what happens when that process changes while the underlying dynamics do not, and we formalise this
\emph{observation-process shift} as a distinct axis of distribution shift. Measuring it across five datasets,
two tasks and two backbones, we find that its severity spans three orders of magnitude and tracks how much
predictive information a task draws from its irregular history: re-observing the test inputs of a wearable-sensor
forecasting benchmark through processes never seen in training raises error by a factor of nineteen, a clinical
forecasting benchmark by a factor of two, and a climate benchmark not at all. The damage is not a function of how
many observations are lost; at matched density, error varies by up to $3.5\times$ depending on \emph{what} is
removed, and the most damaging kind of loss differs by sensing modality. We then develop OPERA, which trains
representations against this shift by exposing the encoder to diverse observation processes, transferring
task-level information from richer observations of the same history, and concentrating training on the
observation processes a model currently handles worst. Exposure to diverse observation processes accounts for
most of the attainable robustness, reducing relative degradation from $+107\%$ to $+14\%$ on the clinical
benchmark and from $+1793\%$ to $+25\%$ on the sensor benchmark at no in-distribution cost, and it transfers to
observation-process \emph{families} absent from training. The other two mechanisms are not interchangeable with
it: privileged-observation distillation improves accuracy where the task is noisy rather than robustness, and
adversarial selection of observation processes helps only when its candidates are harsher than those the model
already trains on. Inference requires neither privileged observations nor candidate search.
\prelim All numbers are single-seed. \todo{finalise with five seeds}
\end{abstract}

\maketitle

\section{Introduction}
Irregular sampling is intrinsic to how many temporal systems are recorded. Laboratory values are ordered when a
clinician judges them necessary; environmental sensors fail, are serviced, and report at heterogeneous rates;
wearable devices lose contact and cycle through power states. In all of these cases the data are not the process
but the outcome of two processes: the temporal dynamics of the system, and the \emph{observation process} that
samples them.

A substantial literature builds encoders for such data, including recurrent models with decay
\cite{che2018grud}, continuous-time latent models \cite{rubanova2019latentode,debrouwer2019gruodebayes,kidger2020ncde},
set-function and attention aggregation \cite{horn2020seft,shukla2021mtand}, graph-guided models
\cite{zhang2022raindrop,zhang2024tpatchgnn}, and query-based input embeddings that let standard multivariate
backbones consume irregular observations without interpolation \cite{lim2026quite}. These methods differ in how
they represent a given set of observations, and they are evaluated under the observation process of the
benchmark on which they were trained.

The observation process, however, is not a fixed property of a domain. A different hospital orders laboratory
panels on a different schedule; a sensor network is redeployed with a different sampling policy; a firmware
update changes when readings are logged. In each case the underlying dynamics are comparable while the
observation process has moved. Our question is whether a representation learned under one observation process
remains useful under another.

We make the distinction explicit (Section~\ref{sec:problem}): the observed sample is $\Oobs=\Sobs(X)$ for an
underlying trajectory $X$ and an observation process $\Sobs \sim P_{\mathrm{obs}}$. \emph{Process shift} changes
$P(X)$; \emph{observation-process shift} changes $P_{\mathrm{obs}}$ while $P(X)$ is comparable. The second is
invisible to a protocol that draws training and test data from a single $P_{\mathrm{obs}}$, which is what
standard benchmarks do.

\paragraph{The phenomenon is real, large, and heterogeneous.}
Measuring it requires only re-observing the \emph{history} of held-out inputs through a different process while
leaving the prediction targets untouched. Doing so across seven task settings (Section~\ref{sec:phenomenon}), we
find mean relative degradations from $+1.9\%$ to $+1793\%$. The ordering is interpretable: severity grows with
how much of a task's predictive signal lives in the irregular history. At one extreme, a climate benchmark
forecasting one month ahead is \emph{improved} by deleting half of its observations, because its signal is
carried by seasonal structure the model can recover without them. At the other, millisecond-resolution
body-sensor dynamics degrade by a factor of thirty-two under regular subsampling alone.

Two further observations shape the method. First, degradation is not a function of observation count: at matched
density, the choice of which observations disappear changes error by up to $3.5\times$. Second, the most damaging
kind of loss is modality-dependent. On clinical data, losing whole variables dominates and temporal
subsampling is nearly free; on body-sensor data the reverse holds. Any mechanism that treats irregularity as a
scalar corruption level cannot express this.

\paragraph{Contributions.}
\begin{itemize}
\item We formalise observation-process shift as a distinct source of distribution shift for irregular time
series, give a composable family of observation processes with validity constraints and a per-sample
descriptor, and define a train/test split over observation processes that makes generalisation to unseen
processes measurable (Sections~\ref{sec:problem}--\ref{sec:setup}).
\item We establish the phenomenon empirically across five datasets, two tasks and two backbones, show that its
severity spans three orders of magnitude and tracks how much a task uses its history, and show that it is not
explained by observation density (Section~\ref{sec:phenomenon}).
\item We develop OPERA, a training framework that acts on the observation process itself through three
mechanisms, and evaluate each in isolation and in combination (Sections~\ref{sec:method}, \ref{sec:results}).
\item We report where each mechanism does and does not help, including two negative results we retain because
they delimit the claim: adversarial observation selection is inert unless its candidates exceed the training
pool in severity, and privileged distillation buys accuracy rather than robustness
(Section~\ref{sec:analysis}).
\end{itemize}

We use a query-based irregular encoder \cite{lim2026quite} as the experimental substrate because it consumes
irregular observations natively and exposes a natural point at which to condition on the observation process.
Nothing in the formulation depends on that choice, and we report results on a second backbone.

\section{Related Work}
\paragraph{Irregular multivariate time series.}
Interpolation-based methods map observations onto a regular grid before encoding \cite{shukla2019ipn};
continuous-time methods evolve a latent state between observations
\cite{rubanova2019latentode,debrouwer2019gruodebayes,kidger2020ncde}; set-based methods aggregate an unordered
observation set with attention \cite{horn2020seft,shukla2021mtand}; graph-based methods model inter-variable
structure \cite{zhang2022raindrop,zhang2024tpatchgnn}. Query-based embeddings \cite{lim2026quite} aggregate the
observations belonging to a variable or temporal patch into a fixed-size token, letting standard backbones
\cite{liu2024itransformer,nie2023patchtst} be reused unchanged. We do not propose a new encoder; we ask how any
such encoder behaves when the observation process moves.

\paragraph{Missingness mechanisms.}
The statistical literature classifies missingness by its dependence on observed and unobserved values
\cite{rubin1976missing}, and in clinical data the missingness pattern is itself informative
\cite{lipton2016informative,che2018grud}. Learning work typically treats irregularity as a property of a dataset
and probes robustness by varying a missingness \emph{rate}. We instead vary the observation \emph{process} at
comparable density and study the case where it differs between training and deployment.

\paragraph{Robustness through exposure and through constraint.}
Randomising nuisance factors during training is standard in sim-to-real transfer
\cite{tobin2017domainrand} and time-series augmentation \cite{wen2021augmentation}; invariant-risk formulations
instead constrain the representation across environments \cite{arjovsky2019irm}. Our observation views randomise
the measurement process rather than the signal. We deliberately do not impose cross-view representation
equality: objectives of that form are known to require explicit anti-collapse terms
\cite{chen2020simclr,grill2020byol,bardes2022vicreg}, and the property we need is that different observation
processes yield representations that remain \emph{useful}, not identical.

\paragraph{Privileged information and distillation.}
Learning with side information available only at training time is the LUPI setting \cite{vapnik2015lupi};
distillation \cite{hinton2015distilling} is one way to realise it, and the two are formally connected
\cite{lopezpaz2016unifying}. A richer observation of the same history is a natural privileged signal for
irregular data, which is the form we adopt.

\paragraph{Hard example and worst-case training.}
Selecting difficult examples accelerates and sharpens training \cite{shrivastava2016ohem}, and training against
a constrained worst case is the basis of adversarial robustness \cite{madry2018towards}. We apply the idea to
the observation process rather than to the input values, which keeps the perturbation set physically meaningful:
every candidate is a valid way the same trajectory could have been recorded.

\section{Learning under Observation-Process Shift}
\label{sec:problem}

\subsection{Setting}
Let $X=\{x(t)\}_{t\in\mathcal{T}}$, $x(t)\in\mathbb{R}^{N}$, be an underlying multivariate temporal process over
a window $\mathcal{T}$. The learner never observes $X$. It receives a finite set of timestamped, partially
observed measurements
\begin{equation}
\Oobs=\{(t_k,\,x(t_k),\,m_k)\}_{k=1}^{K},\qquad m_k\in\{0,1\}^{N},
\end{equation}
where $m_{k,i}=1$ iff variable $i$ was measured at time $t_k$. We call the mechanism that selects
$\{(t_k,m_k)\}$ an \emph{observation process} $\Sobs$, write $\Oobs=\Sobs(X)$, and model it as a random variable
$\Sobs\sim P_{\mathrm{obs}}$. A model $f_\theta$ maps $\Oobs$ to a task output $\hat{Y}$: a forecast of $x$ over
a future window, or a label.

Training data are drawn from a joint distribution over $(X,\Sobs)$. We distinguish
\begin{align}
\text{process shift:}\quad & P_{\mathrm{train}}(X)\ \neq\ P_{\mathrm{test}}(X), \\
\text{observation-process shift:}\quad & P_{\mathrm{train}}(\Sobs)\ \neq\ P_{\mathrm{test}}(\Sobs),
\end{align}
with $P(X)$ comparable in the second case. This paper concerns the second. Two datasets may record the same
underlying process while differing in observation density, in the temporal clustering of observations, in which
variables are observed, or in sampling regularity.

\subsection{A composable family of observation processes}
\label{sec:family}
We need observation processes that are \emph{realisable}: each must correspond to a way the same trajectory
could have been recorded, so that the target remains well defined. We therefore restrict attention to processes
that act on the observation mask by \emph{removal} only. Formally, an observation process is a map
\begin{equation}
\Sobs:(t,m)\ \mapsto\ (t,\,m\odot\kappa),\qquad \kappa\in\{0,1\}^{K\times N},
\end{equation}
and values at removed entries are discarded rather than imputed, so $\Sobs(X)$ is always a sub-sample of the
original record and no artificial value is ever introduced. This is what allows a single held-out set to be
re-observed many times with the supervision untouched.

We use nine primitive families, grouped by the structure they destroy:
\begin{itemize}
\item \textbf{unstructured}: \textsc{RandomDropout}$(p)$ removes each observation independently with
probability $p$; \textsc{DensityControlledSampling}$(\rho)$ retains a fixed fraction $\rho$ of every variable's
observations, fixing density independently of the original pattern.
\item \textbf{temporal}: \textsc{TemporalStride}$(k)$ keeps every $k$-th observation of each variable, lowering
the sampling \emph{rate} while keeping it regular; \textsc{BurstMissingness}$(p)$ removes contiguous windows
covering a fraction $p$ of the observed span; \textsc{LongGapMissingness}$(p)$ is the single-window case;
\textsc{LocalTemporalMask}$(p)$ removes one contiguous run per variable, matching random dropout in density but
clustering the loss per channel.
\item \textbf{variable-wise}: \textsc{VariableDropout}$(p)$ removes whole variables independently;
\textsc{VariableGroupDropout}$(p)$ removes contiguous blocks of variables together, modelling a failed
monitor or an unavailable panel.
\end{itemize}
Processes compose sequentially ($\Sobs_2\circ\Sobs_1$, written \texttt{a+b}) and mix
($\Sobs\sim\mathrm{Unif}\{\Sobs_1,\dots\}$, written \texttt{a|b}), so an observation regime is fully described by
a string and can be specified in configuration rather than code. \textsc{LocalTemporalMask} and
\textsc{DensityControlledSampling} exist specifically to construct density-matched comparisons
(Section~\ref{sec:phenomenon}).

\subsection{Descriptor and validity constraints}
\label{sec:descriptor}
Two auxiliary objects are needed by the mechanisms. First, a per-sample \emph{descriptor}
$r=g(\Oobs)\in\mathbb{R}^{8}$ summarising the realised observation process: observation count and density, the
mean, standard deviation and maximum of inter-observation gaps, temporal coverage, variable coverage, and recent
density. It is computed from the mask and timestamps alone, never from values, so it can be logged, used to
select processes, or fed to the encoder without leaking supervision.

Second, a set of \emph{validity constraints} $\mathcal{C}$, since a mechanism that searches over observation
processes must be prevented from selecting degenerate ones. We constrain relative density, variable coverage,
maximum gap, per-variable coverage and minimum observation count:
\begin{equation}
\mathcal{C}:\quad \rho_{\min}\le\frac{\rho(\Sobs(X))}{\rho(X)}\le\rho_{\max},\quad
\Delta t_{\max}(\Sobs(X))\le G_{\max},\quad \dots
\label{eq:constraints}
\end{equation}
A candidate violating any bound is discarded before selection. Throughout we use $\rho_{\min}=0.15$, variable
coverage $\ge 0.3$, $G_{\max}=0.6$ of the window, and at least four observations per sample.

\section{OPERA}
\label{sec:method}
OPERA acts on $\Sobs$, not on the encoder. All three mechanisms manipulate the same object --- the observation
process through which the model sees a trajectory --- and they differ in \emph{how} they choose it:
augmentation samples it, distillation contrasts a sparse choice against a richer one, and mining searches for
the worst admissible choice. The encoder, its parameter count and its inference path are untouched.

\subsection{Mechanism 1: observation-process augmentation}
\label{sec:aug}
The first mechanism treats the observation process as nuisance variation to be randomised over. For each
training trajectory we draw $V$ processes from a mixture $P_{\mathrm{mix}}$ and train on all of the resulting
views:
\begin{equation}
\Oobs_v=\Sobs_v(X),\quad \Sobs_v\sim P_{\mathrm{mix}},\qquad
\mathcal{L}_{\mathrm{task}}=\frac1V\sum_{v=1}^{V}\ell\big(f_\theta(\Oobs_v),\,Y\big),
\label{eq:aug}
\end{equation}
with $\ell$ the masked mean squared error for forecasting and cross-entropy for classification. The targets $Y$
are shared across views and never modified: every view is a different recording of the same episode.

We make no representation-level requirement. The intended property is that the representation of a trajectory
stay \emph{task-useful} under a change of observation process, which is weaker than
$z(\Sobs_1(X))=z(\Sobs_2(X))$ and, unlike it, does not admit the degenerate solution of collapsing all
representations together.

\subsection{Mechanism 2: privileged-observation distillation}
\label{sec:distill}
A sparse observation may simply not determine the target. During training, however, we can often construct a
\emph{richer observation of the same history}. This is privileged information in the sense of
\cite{vapnik2015lupi}: available at training time, absent at deployment.

Let $\Sobs_{\mathrm{rich}}$ be a process at least as dense as the student's view, in our experiments the
untouched record. A teacher $f_{\theta_T}$ is trained on rich observations and frozen; the student is trained on
sparse views and additionally imitates the teacher's \emph{task output}:
\begin{equation}
\mathcal{L}_{\mathrm{distill}}
=\big\| f_\theta(\Sobs_v(X))-\mathrm{sg}\big[f_{\theta_T}(\Sobs_{\mathrm{rich}}(X))\big] \big\|^2_{\Omega},
\label{eq:distill}
\end{equation}
where $\mathrm{sg}[\cdot]$ is the stop-gradient and $\Omega$ restricts the norm to entries carrying supervision,
so the student is never asked to imitate the teacher where no target exists. For probabilistic outputs we also
support a KL form. We additionally implement an exponential-moving-average teacher,
$\theta_T\leftarrow\alpha\theta_T+(1-\alpha)\theta$, and treat the choice as an empirical question.

Distillation is deliberately at the level of predictions rather than representations. A representation-matching
variant exists in our implementation but is reported only as an ablation: it reintroduces exactly the
constraint that Section~\ref{sec:aug} declines to impose.

\paragraph{Anti-leakage invariant.}
Teacher and student must differ \emph{only} in the observation of the history. Both are evaluated on the same
history window, the same prediction times and the same supervision; the teacher never sees a future value or a
label the student cannot use. We enforce this with an assertion on every batch key that carries supervision and
verify it in unit tests, for both forecasting and classification.

\subsection{Mechanism 3: hard observation-process mining}
\label{sec:mining}
Random exposure spends most of its budget on observation processes the model already handles. The third
mechanism instead searches for the admissible process that the \emph{current} model handles worst,
\begin{equation}
\Sobs^{*}=\arg\max_{\Sobs\in\mathcal{A}}\ \ell\big(f_\theta(\Sobs(X)),Y\big),
\qquad \mathcal{A}=\{\Sobs: \Sobs \text{ satisfies } \mathcal{C}\},
\label{eq:mining}
\end{equation}
with $\mathcal{C}$ from Eq.~\ref{eq:constraints}. We approximate the maximisation by scoring a finite candidate
pool rather than differentiating through mask generation: $K$ candidates are drawn from a pool $\mathcal{P}$,
evaluated under \texttt{no\_grad}, filtered through $\mathcal{C}$, and the hardest retained. Only the selected
views enter the backward pass, so mining costs $K$ extra forward passes and no extra gradient memory.

Two design points proved to matter. First, training exclusively on hard views destabilises optimisation and
biases the model away from the deployment regime, so we mix,
\begin{equation}
P(\Sobs)=(1-\alpha)\,P_{\mathrm{mix}}+\alpha\,P_{\mathrm{hard}},
\label{eq:alpha}
\end{equation}
and additionally hold $\alpha=0$ for a warm-up period, giving the curriculum: random exposure, then mixed random
and mined. Second, and more consequentially, \emph{the candidate pool must be harsher than the training pool}.
If $\mathcal{P}=P_{\mathrm{mix}}$ then the hardest of $K$ draws is close in severity to a random draw and mining
degenerates to augmentation at $K$ times the cost. We therefore treat $\mathcal{P}$ as a separate
hyper-parameter and report both choices (Section~\ref{sec:results}).

\subsection{Objective and training procedure}
The complete objective adds one term to the task loss,
\begin{equation}
\boxed{\ \mathcal{L}=\mathcal{L}_{\mathrm{task}}+\lambda_{d}\,\mathcal{L}_{\mathrm{distill}}\ }
\label{eq:objective}
\end{equation}
where mining determines \emph{which} observation processes contribute to $\mathcal{L}_{\mathrm{task}}$ in
Eq.~\ref{eq:aug}. No further loss is introduced, and in particular no representation-consistency term. Setting
$\lambda_d=0$, $\alpha=0$ and $V=1$ recovers ordinary training.

\begin{algorithm}[t]
\caption{OPERA training step}
\label{alg:opera}
\begin{algorithmic}[1]
\Require batch $(X,Y)$ with observation $(t,m)$; epoch $e$; mixture $P_{\mathrm{mix}}$; pool $\mathcal{P}$;
constraints $\mathcal{C}$; frozen teacher $f_{\theta_T}$
\If{$e<e_{\mathrm{warm}}$ \textbf{or} $\alpha=0$}
  \State $\{\Sobs_v\}_{v=1}^{V}\sim P_{\mathrm{mix}}$ \Comment{stage 1: random exposure}
\Else
  \State draw $K$ candidates from $\mathcal{P}$; keep those satisfying $\mathcal{C}$
  \State score survivors with $f_\theta$ under \texttt{no\_grad}; take the $\lceil\alpha V\rceil$ hardest
  \State fill the remaining $V-\lceil\alpha V\rceil$ views from $P_{\mathrm{mix}}$
\EndIf
\State $\hat{Y}_T \gets \mathrm{sg}\big[f_{\theta_T}(\Sobs_{\mathrm{rich}}(X))\big]$ \Comment{same history}
\State $\mathcal{L}\gets\frac1V\sum_v \ell(f_\theta(\Sobs_v(X)),Y) + \lambda_d\sum_v \mathcal{L}_{\mathrm{distill}}$
\State one optimiser step on $\mathcal{L}$
\end{algorithmic}
\end{algorithm}

Algorithm~\ref{alg:opera} is a single unified update: one loss, one backward pass, one optimiser step. The only
staging is the mining warm-up and the fact that a frozen teacher is trained beforehand.

\paragraph{Deployment.}
At inference the model consumes one sparse observation and emits a prediction. The teacher, the rich view and
the candidate pool are training-time constructs; parameter count and inference latency are identical to the
unmodified encoder (Table~\ref{tab:cost}).

\paragraph{Optional observation-aware encoding.}
The descriptor $r=g(\Oobs)$ of Section~\ref{sec:descriptor} can also be supplied to the encoder, conditioning
aggregation on the observation regime. We implement this but keep it outside the framework proper: the
mechanisms above require no architectural change, and we prefer to establish them first. \todo{report the
observation-aware ablation}

\section{Experimental Setup}
\label{sec:setup}

\subsection{Datasets and tasks}
We use five datasets spanning three sensing modalities, two tasks and two backbones
(Table~\ref{tab:datasets}). Preprocessing, splits and normalisation are the unmodified published pipelines:
t-PatchGNN \cite{zhang2024tpatchgnn} for the forecasting benchmarks (PhysioNet from the 2012 challenge
\cite{silva2012physionet}) and Raindrop \cite{zhang2022raindrop} for classification, including the minority-class upsampling of the binary tasks. USHCN appears at two horizons
because, as Section~\ref{sec:phenomenon} shows, horizon length changes how much a model must rely on its
history.

\begin{table}[t]
\caption{Datasets and task settings. Backbone follows the published protocol per task.}
\label{tab:datasets}
\small
\begin{tabular}{llrrll}
\toprule
setting & task & vars & samples & backbone & metric \\
\midrule
PhysioNet 24h$\to$24h & forecast & 41 & 12{,}000 & iTransformer & MSE \\
Activity 3000$\to$1000\,ms & forecast & 12 & 5{,}400 & iTransformer & MSE \\
USHCN 24$\to$1 mo & forecast & 5 & 26{,}736 & iTransformer & MSE \\
USHCN 24$\to$12 mo & forecast & 5 & 26{,}736 & iTransformer & MSE \\
P12 mortality & classify (2) & 36 & 11{,}988 & PatchTST & AUROC \\
P19 sepsis & classify (2) & 34 & 38{,}803 & PatchTST & AUROC \\
PAM activity & classify (8) & 17 & 5{,}333 & PatchTST & accuracy \\
\bottomrule
\end{tabular}
\end{table}

\subsection{Observation-process splits}
\label{sec:splits}
The central protocol decision is that training and evaluation draw from \emph{disjoint} observation-process
distributions, so that generalisation across observation processes is measurable rather than assumed.

\textbf{Training pool} $P_{\mathrm{mix}}$ is deliberately mild: random dropout at $p\in\{0.1,0.25,0.4\}$,
stride $2$, burst at $p\in\{0.2,0.3\}$, variable dropout at $p=0.2$, and local masking at $p=0.3$.

\textbf{Evaluation conditions} comprise 16 processes: the training families at unseen severities, plus families
absent from training. Nine are \emph{unseen}: random dropout $0.7$, variable dropout $0.5$, stride $3$ and $4$,
long gap $0.4$, burst $0.5$, variable-group dropout $0.3$, and two compositions (random $0.5$ $+$ variable
$0.3$; stride $3$ $+$ variable-group $0.4$). We report degradation over all shifted conditions and separately
over the unseen subset.

\textbf{Evaluation procedure.} For each condition we re-observe the \emph{history} of every test input and leave
the prediction targets and labels untouched. The pseudo-random draw is seeded identically for every model, so
all models are compared on the same re-observed inputs. The clean condition is the dataset's own observation
process.

\subsection{Metrics}
Forecasting uses masked MSE and MAE over observed target entries, averaged per variable then over variables with
supervision. Binary classification uses AUROC and AUPRC; the eight-class task uses accuracy and macro F1. We
summarise sensitivity by relative degradation with respect to a model's own clean score,
\begin{equation}
\Delta_{\mathrm{shift}}=\frac{1}{|\mathcal{E}|}\sum_{\Sobs\in\mathcal{E}}
\frac{\mathrm{err}_{\Sobs}-\mathrm{err}_{\mathrm{clean}}}{\mathrm{err}_{\mathrm{clean}}},
\label{eq:degradation}
\end{equation}
over the shifted conditions $\mathcal{E}$, and report $\Delta_{\mathrm{shift}}$ restricted to the unseen subset
as well. $\Delta_{\mathrm{shift}}$ must be read \emph{beside} absolute error and never instead of it: because it
is normalised by a model's own clean score, a uniformly weaker model can attain a flattering ratio. We therefore
report clean, mean shifted and worst-case error in every table. We also report parameter count, training time
per epoch, peak training memory and inference latency.

\subsection{Protocol}
All models use the query-based irregular embedding with hidden width 64, four heads and three layers, trained
with Adam at $10^{-3}$, batch 32 (128 for USHCN; 16 for P12, where a single view's attention already occupies
4.7\,GB), an epoch cap of 40, early-stopping patience 10 and a minimum of 15 epochs. Runs are executed
sequentially on one 8\,GB GPU. Every run records its configuration, git commit, seed, checkpoint and metrics.
\textbf{All reported numbers use a single seed} and are marked \prelim; the protocol calls for five, and we
regard nothing here as a statistical result. \todo{five seeds for the primary comparisons}

\section{Does Observation-Process Shift Matter?}
\label{sec:phenomenon}
We first establish the phenomenon on models trained without any OPERA mechanism.

\paragraph{Severity spans three orders of magnitude.}
Table~\ref{tab:phenomenon} reports each baseline under its own observation process and under shift. Mean
degradation ranges from $+1.9\%$ to $+1793\%$, and the ordering is interpretable: it grows with how much of the
task's predictive signal must come from the irregular history. USHCN forecasting one month ahead from five
climate variables sits at the bottom, and the diagnostic is unambiguous --- removing half of its observations at
matched density \emph{lowers} its error, from $0.5379$ to $0.5355$. Its signal is carried by seasonal structure
that survives the loss. Extending the same dataset to a twelve-month horizon doubles its sensitivity
($+1.9\%\to+3.7\%$), which identifies the mechanism of the null rather than excusing it. At the top, Activity's
millisecond-resolution sensor dynamics degrade by $32\times$ under stride-4 subsampling alone.

\begin{table}[t]
\caption{Observation-process shift on models trained without any OPERA mechanism\prelim. Primary metric per
setting; degradation by Eq.~\ref{eq:degradation}. Every shifted condition re-observes only the history.}
\label{tab:phenomenon}
\small
\begin{tabular}{lrrrr}
\toprule
setting & clean & worst condition & $\Delta_{\mathrm{shift}}$ & $\Delta_{\mathrm{unseen}}$ \\
\midrule
Activity (MSE) & 0.00263 & 0.08547 \textit{(stride 4)} & $+1793\%$ & $+2057\%$ \\
PhysioNet (MSE) & 0.00549 & 0.02134 \textit{(var .5)} & $+107\%$ & $+130\%$ \\
PAM (acc.) & 88.01 & 53.93 \textit{(var .5)} & $+64\%$ & $+88\%$ \\
USHCN-12mo (MSE) & 0.52020 & 0.55338 \textit{(comp.)} & $+3.7\%$ & $+4.6\%$ \\
P19 (AUROC) & 82.13 & 74.63 \textit{(var .3)} & $+2.9\%$ & $+3.3\%$ \\
USHCN-1mo (MSE) & 0.53787 & 0.56933 \textit{(var .5)} & $+1.9\%$ & $+2.8\%$ \\
P12 (AUROC) & 83.53 & 77.98 \textit{(comp.)} & \multicolumn{2}{c}{\textit{see \S\ref{sec:analysis}}} \\
\bottomrule
\end{tabular}
\end{table}

\paragraph{Density does not explain it.}
Table~\ref{tab:density} compares processes that remove approximately half of the observations but destroy
different structure. On PhysioNet the spread is $2.9\times$ and on PAM $3.5\times$ in cross-entropy; a model can
lose half its observations almost for free if they are removed uniformly in time, and catastrophically if they
are removed by variable. Treating irregularity as a scalar corruption level cannot represent this distinction,
which is the empirical case for modelling the observation process as an object.

\begin{table}[t]
\caption{Matched-density comparison\prelim: four processes each removing $\approx$50\% of observations, on
models trained without OPERA. Error varies by $2.9$--$3.5\times$ with \emph{what} is removed.}
\label{tab:density}
\small
\begin{tabular}{lrrrr}
\toprule
 & density- & stride & burst & variable \\
setting & matched & 2 & .5 & dropout .5 \\
\midrule
PhysioNet (MSE) & 0.00736 & 0.00791 & 0.00804 & \textbf{0.02134} \\
Activity (MSE) & 0.04892 & 0.05230 & \textbf{0.01973} & 0.05778 \\
PAM (cross-ent.) & 0.4466 & 0.4360 & 0.4897 & \textbf{1.5578} \\
\bottomrule
\end{tabular}
\end{table}

\paragraph{The damaging kind of loss is modality-dependent.}
On the clinical datasets, losing whole variables dominates while temporal subsampling is nearly free: on P19,
removing 30\% of variables costs 7.5 AUROC whereas stride 3 costs 1.1; on PAM, removing half the variables costs
34 accuracy points whereas stride 3 costs 1.5. Activity inverts this, with stride its single worst condition and
burst removal its mildest (Table~\ref{tab:density}). Hourly laboratory values are largely redundant in time and
irreplaceable across channels; body-sensor dynamics are the reverse. A robustness method must therefore cover
process \emph{families}, not a severity axis.

\section{Results}
\label{sec:results}

\paragraph{Exposure to diverse observation processes accounts for most of the attainable robustness.}
Table~\ref{tab:components} isolates the mechanisms on PhysioNet. Augmentation alone reduces
$\Delta_{\mathrm{shift}}$ from $+107\%$ to $+14\%$ and $\Delta_{\mathrm{unseen}}$ from $+130\%$ to $+18\%$ while
leaving clean error unchanged ($0.00553$ vs.\ $0.00549$), and it improves all sixteen conditions. It transfers to
process families absent from training: stride subsampling, never trained on in any form, falls from $0.00923$ to
$0.00572$. On Activity the same mechanism, combined with the others, reduces mean shifted error from $0.04983$ to
$0.00358$, a factor of fourteen (Table~\ref{tab:main}).

\begin{table}[t]
\caption{Mechanism isolation on PhysioNet (test MSE\prelim). ``pool'' is the hard-mining candidate pool of
Section~\ref{sec:mining}. Read $\Delta$ beside the absolute columns.}
\label{tab:components}
\small
\begin{tabular}{lrrrr}
\toprule
configuration & clean & mean shift & worst & $\Delta_{\mathrm{shift}}$ \\
\midrule
no mechanism & \textbf{0.00549} & 0.01137 & 0.02134 & $+107.0\%$ \\
augmentation & 0.00553 & 0.00632 & 0.00772 & $+14.3\%$ \\
\quad $+$ distillation & 0.00573 & 0.00644 & 0.00776 & $+12.3\%$ \\
\quad $+$ mining, pool $=P_{\mathrm{mix}}$ & 0.00558 & 0.00641 & 0.00765 & $+14.8\%$ \\
\quad $+$ mining, harsher pool & 0.00562 & \textbf{0.00626} & \textbf{0.00721} & $+11.4\%$ \\
all three, pool $=P_{\mathrm{mix}}$ & 0.00559 & 0.00632 & 0.00741 & $+13.1\%$ \\
all three, harsher pool & 0.00590 & 0.00653 & 0.00748 & $\mathbf{+10.7\%}$ \\
\bottomrule
\end{tabular}
\end{table}

\paragraph{Adversarial selection of observation processes helps only when its candidates exceed the training
pool.} With $\mathcal{P}=P_{\mathrm{mix}}$, mining is indistinguishable from augmentation ($+14.8\%$ vs.\
$+14.3\%$; mean shifted $0.00641$ vs.\ $0.00632$) at $2.7\times$ the training cost, because the hardest of $K$
draws from the training mixture is barely harder than one draw. Given a pool containing severities and
compositions absent from training, the same mechanism becomes the strongest single configuration on absolute
shifted error and worst case ($0.00626$ and $0.00721$), with the gain concentrated exactly where it should be:
$0.00706$ vs.\ $0.00747$ under 70\% dropout and $0.00721$ vs.\ $0.00753$ under 50\% variable removal. This is a
property of the pool, not of the mechanism, and it is easy to miss: the natural default --- mining over the
processes one already trains on --- yields a null result.

\paragraph{Privileged distillation improves accuracy where the task is noisy, not robustness.}
On PhysioNet it is within noise of augmentation. On USHCN, where the task is hard and the student underfits, it
is the largest single effect we measure: clean MSE $0.5379\to0.5135$ ($-4.5\%$), best on thirteen of sixteen
conditions, while $\Delta_{\mathrm{shift}}$ does not improve. On P19 the full framework lifts clean AUROC
$82.13\to85.59$ and AUPRC $16.89\to24.47$ ($+45\%$). The pattern is consistent with distillation acting as a
smoother regression target rather than as a robustness mechanism.

\paragraph{Across settings, benefit tracks baseline damage.}
Table~\ref{tab:main} compares no mechanism against the full framework on every setting. Where degradation is
large the gains are large; where the baseline is already insensitive there is little to recover, and the
framework neither helps nor harms materially. On the classification tasks the effect appears in the primary
metric: PAM recovers $10.5$ and $11.1$ accuracy points at 30\% and 50\% variable removal, P19 gains $5.6$ AUROC
at 30\% variable removal and improves clean accuracy as well, and P12 --- which barely degrades --- gains $2.6$
AUROC at 50\% variable removal for a $0.2$ AUROC clean cost.

\begin{table}[t]
\caption{No mechanism vs.\ the full framework across all settings\prelim. Primary metric; MSE lower-better,
AUROC/accuracy higher-better. $\Delta$ is Eq.~\ref{eq:degradation} on the primary error.}
\label{tab:main}
\small
\begin{tabular}{llrrrr}
\toprule
 &  & \multicolumn{2}{c}{no mechanism} & \multicolumn{2}{c}{full OPERA} \\
\cmidrule(lr){3-4}\cmidrule(lr){5-6}
setting & metric & clean & $\Delta$ & clean & $\Delta$ \\
\midrule
Activity & MSE & \textbf{0.00263} & $+1793\%$ & 0.00286 & $\mathbf{+25\%}$ \\
PhysioNet & MSE & \textbf{0.00549} & $+107\%$ & 0.00559 & $\mathbf{+13\%}$ \\
PAM & acc. & \textbf{88.01} & $+64\%$ & 87.27 & $\mathbf{+34\%}$ \\
USHCN-12mo & MSE & 0.52020 & $+3.7\%$ & \textbf{0.51686} & $\mathbf{+2.6\%}$ \\
P19 & AUROC & 82.13 & $+2.9\%$ & \textbf{85.59} & $\mathbf{+1.2\%}$ \\
USHCN-1mo & MSE & 0.53787 & $+1.9\%$ & \textbf{0.52107} & $\mathbf{+1.5\%}$ \\
P12 & AUROC & \textbf{83.53} & --- & 83.30 & --- \\
\bottomrule
\end{tabular}
\end{table}

\paragraph{Cost.}
Table~\ref{tab:cost} reports the overhead. Every mechanism is training-time only: parameter count is identical
across configurations and inference latency is unchanged at $0.63$\,ms per series. Augmentation costs
$1.7\times$ time and $1.7\times$ peak memory for two views; mining adds $K$ forward passes per step.

\begin{table}[t]
\caption{Training and deployment cost on PhysioNet\prelim. Inference and parameter count are unaffected by
every mechanism.}
\label{tab:cost}
\small
\begin{tabular}{lrrrr}
\toprule
configuration & params & s/epoch & peak GB & infer.\ ms \\
\midrule
no mechanism & 129{,}153 & 24.1 & 3.7 & 0.630 \\
augmentation & 129{,}153 & 41.4 & 6.3 & 0.670 \\
$+$ distillation & 129{,}153 & 45.3 & 6.3 & 0.630 \\
$+$ mining (harsher pool) & 129{,}153 & 82.3 & 6.3 & 0.634 \\
all three (harsher pool) & 129{,}153 & 86.3 & 6.3 & 0.629 \\
\bottomrule
\end{tabular}
\end{table}

\section{Analysis and Limitations}
\label{sec:analysis}

\paragraph{What the negative and partial results delimit.}
Two mechanisms did not behave as their motivation suggested, and both boundaries are informative. Mining's
inertness under its natural default shows that "train on the worst case" is empty unless the search space
exceeds what random exposure already covers. Distillation's profile shows that a richer observation helps when
the student cannot fit the target, which is a statement about task noise rather than about observation-process
shift; the two goals are separable and should not be conflated because both happen to involve sparse views.
Conversely, augmentation's dominance is itself a finding: informing a model by \emph{showing} it varied
observation processes was consistently more reliable than constraining what its representation must look like.

\paragraph{When the ratio metric misleads.}
P12's cross-entropy \emph{improves} under mild shifts while its AUROC and AUPRC fall, so
$\Delta_{\mathrm{shift}}$ computed on cross-entropy is negative and meaningless there; we omit it in
Table~\ref{tab:main} and read P12 from the ranking metrics. More generally, because
Eq.~\ref{eq:degradation} normalises by a model's own clean score, a uniformly weaker model can post a better
ratio. On USHCN-1mo the baseline's ratio ($+1.9\%$) is close to the framework's ($+1.5\%$) while its shifted
error is higher everywhere, and its clean error exceeds the framework's error \emph{under shift}. Absolute
error remains the quantity of interest. \todo{recompute degradation on AUROC for binary tasks}

\paragraph{Limitations.}
Every number is single-seed; the protocol requires five, and the smaller differences in
Tables~\ref{tab:components} and \ref{tab:main} are within the range we would expect seeds to move. Observation
processes are synthetic re-observations of real records, which is what makes the targets well defined but is
weaker evidence than a shift between two real acquisition environments. Compositions are limited to pairs. The
two backbones are tied to tasks rather than crossed with them. USHCN is not min-max normalised by its published
pipeline, so its errors are in raw units and not comparable in magnitude to the other forecasting settings. The
observation-aware encoder of Section~\ref{sec:method} is implemented but not evaluated here.

\section{Conclusion}
Irregularly sampled data should be treated as the product of a temporal process and an observation process, and
the second can shift independently of the first. We formalised that shift, showed that its severity varies by
three orders of magnitude across benchmarks in a way that tracks how much each task uses its irregular history,
and showed that it is not reducible to how many observations are lost. Of the mechanisms we studied for acting
on the observation process, randomising it during training is the most reliable and transfers to process
families never seen; transferring task-level information from a richer observation of the same history improves
accuracy where the task is noisy; and searching for worst-case observation processes helps only when the search
space is harsher than what random exposure already supplies. None of them changes what deployment requires: one
sparse observation and the student.

\bibliographystyle{ACM-Reference-Format}
\bibliography{refs}

\end{document}
